\section{Literature review}
\label{sec:lit}

\paragraph{Hydropower-dependent systems and climate risk.}
Hydro-dominated systems face supply risk that is driven by water availability rather than by generating
capacity. At the global scale, van Vliet et al.\ \cite{vanvliet2016vulnerability} showed that a large share of
hydropower and thermoelectric capacity is vulnerable to reduced streamflow and warmer water under climate change.
They also argued that adaptation options such as efficiency gains and fuel switching matter. For Albania,
IRENA's readiness assessment highlights both the dominance of the Drin cascade and the projected decline in
hydro output \cite{irena2021rra}. Reservoir operation has a long literature of its own
\cite{labadie2004reservoir}. Stochastic dual dynamic programming (SDDP) \cite{pereira1991sddp} remains the
reference method for multi-stage hydrothermal scheduling under inflow uncertainty. We use a simpler two-stage
structure that is adequate for investment screening (Section~\ref{sec:formulation}), and we note SDDP as a
natural extension.

\paragraph{Renewable capacity expansion and open models.}
Linear capacity-expansion models that co-optimise investment and hourly dispatch are standard tools
\cite{pfenninger2014modeling}. PyPSA \cite{brown2018pypsa} and its European application PyPSA-Eur
\cite{horsch2018pypsaeur} established open, reproducible practice for such models. Studies with them have
shown, for example, the value of transmission cooperation in highly renewable systems
\cite{schlachtberger2017cooperation}. Weather-derived capacity factors from reanalysis are widely used when
measured output is unavailable. They require care, however, because reanalysis wind speeds are biased and
need correction at the site level \cite{staffell2016wind,pfenninger2016pv,gelaro2017merra2,hersbach2020era5}.

\paragraph{Flexibility and storage.}
Reviews of flexibility options \cite{lund2015flexibility,kondziella2016flexibility} distinguish short-duration
resources (batteries, demand response) from long-duration ones (reservoir hydro, interconnection). The value
of short-duration storage depends strongly on the temporal resolution of the model: coarse time steps average
away the intra-day variability that storage monetises \cite{poncelet2016temporal}. Representative-period
methods reduce computation but must treat seasonal storage carefully \cite{kotzur2018seasonal}. Projected
storage costs continue to fall \cite{schmidt2019storage}. The market value of variable renewables declines as
their penetration rises \cite{hirth2013market}, which is directly relevant to solar-heavy expansion in a
small system.

\paragraph{Stochastic and risk-sensitive planning.}
Two-stage stochastic programming \cite{birge2011stochastic} separates here-and-now investment from
wait-and-see operation. The value of the stochastic solution (VSS) and the expected value of perfect
information (EVPI) quantify, respectively, the benefit of modelling uncertainty and the cost of not knowing
the future \cite{birge2011stochastic,madansky1960inequalities}. Conditional value at risk (CVaR) admits a
linear representation \cite{rockafellar2000cvar,rockafellar2002cvar} and is widely used in electricity
decision problems \cite{conejo2010decision}. Reliability is commonly priced through a value of lost load
(VOLL), whose empirical estimates vary by an order of magnitude across countries and customer groups
\cite{schroder2015voll}. This motivates the sensitivity analysis reported below.

\paragraph{Gap.}
We are not aware of an open, reproducible model of Albania that combines (i)~a physically explicit reservoir
water balance, (ii)~independent battery power and energy investment, (iii)~an interconnector that is limited
rather than infinite, (iv)~two-stage stochastic investment with CVaR, and (v)~a marginal-value analysis of
flexibility with dual--finite-difference reconciliation. This work offers such a model as a transparent
research prototype.
